\scanpage{174}
% L1–8
peut être interprétée, par l'identification
\[
\begin{aligned}
{}^\lambda A&\xrightarrow{\sim}A^0\\
x&\mapsto x-t
\end{aligned}
\]
comme une application de $A$ dans $A^0$,
de façon précise on définit
\[
u:A\longrightarrow A^0\qquad u(x)=\lambda x-t
\]
application affine compatible avec
\[
A^0\mathrel{\mathop{\text{---}}\limits^{\lambda\,\mathrm{id}_{A^0}}}A^0
\]
% L9–10, between lines
\inserted{et inversement, une telle application s'interprète comme
\deleted{\unsure{application|applic. en}} \unsure{isomorphisme|iso} de $A^0$-torseurs $ {}^\lambda A\simeq A^0$, i.e. choix d'une origine \unsure{$*$}}
\margin{$*$ de $ {}^\lambda A$}

% L10–15
Dans le cas où $\lambda=2$, donc \unsure{où $t$ correspond}
à une involution de $F$ au sens
du début de ces notes, on a donc
\[
u(x)=2x-t=-(\sigma(x)-x),
\]
c'est donc au signe près notre fameux
$\sigma(x)-x$.

% L16–22
Or dans notre construction \unsure{de|des}
$*A_i$, ce sont ces applications $u_i(x_i)$
(ou ce qui revient au $=$, au signe
près, les $\sigma_i(x_i)-x_i$) qui jouaient le
rôle crucial. \deleted{Il n'y avait pas}
Cette application $u(x)$ étant affine,
se prolonge linéairement en
\[
E\xrightarrow{u}A^0\qquad
u(x)\unsure{\mathrel{=\!\!-}|=}\lambda x-\underbrace{\epsilon(x)t}_{=\pi\epsilon(x)}
\]
% archive sheet number: 173
